% This is a variation of the NeurIPS'24 format.
\documentclass{article}
\usepackage[final]{adrl}

\usepackage[utf8]{inputenc}
\usepackage[T1]{fontenc}
\usepackage{hyperref}
\usepackage{url}
\usepackage{booktabs}
\usepackage{amsfonts}
\usepackage{amsmath}
\usepackage{nicefrac}
\usepackage{microtype}
\usepackage{xcolor}

\title{Can Intrinsic Rewards Ignore the Noisy TV?\\
A MiniGrid Study of Reducible vs.\ Irreducible Novelty}
\author{Christian Jesinghaus\\
\url{https://github.com/ChristianJesinghaus/RL-exercises}}

\begin{document}
\maketitle

\section{Motivation}
Intrinsic rewards can help exploration under sparse extrinsic rewards, but
prediction-error bonuses can fail on irrelevant stochastic observations. This
failure is known as the \emph{noisy-TV problem}: random observations remain
surprising although they do not help solve the task. I study whether a simple
Random Network Distillation (RND) variant based on learning progress reduces
this failure.

\textbf{Research question:}
\begin{quote}
Can learning-progress-gated RND avoid noisy-TV distraction while preserving
the exploration benefit of intrinsic rewards in sparse-reward MiniGrid?
\end{quote}

Here, \emph{robustness} means maintaining extrinsic task performance without
spending excessive time at an irrelevant stochastic distractor.

\section{Related Topics}
The project builds on Deep Q-Networks (DQN)~\citep{mnih-nature15},
Double DQN~\citep{vanhasselt-aaai16}, count-based exploration as
background~\citep{bellemare-nips16}, and RND~\citep{burda-iclr19}.
Kayal et al.~\citep{kayal2025impact} compare intrinsic rewards in MiniGrid using
return and coverage metrics. The contribution here is a controlled noisy-TV
stress test and a learning-progress gate inspired by work on irreducible
uncertainty and noisy-TV avoidance~\citep{mavorparker2022noisytv,hou2025beyond}.

\section{Method}
\paragraph{Environment.}
I use a controlled $11\times11$ FourRooms variant implemented with
MiniGrid~\citep{minigrid2023}. It has a fixed start, goal, and four door
openings; the route to the goal crosses two room boundaries. The compact layout
keeps walls and openings locally identifiable to a feed-forward DQN while
retaining MiniGrid's partial $7\times7\times3$ observation.

A fixed region near the start acts as the noisy TV. While the agent is in this
region, the environment appends a randomly resampled, task-irrelevant vector to
the partial observation. The distractor changes neither transitions, extrinsic
reward, goal position, nor the shortest path. All compared methods receive the
same full observation: symbolic image, MiniGrid's direction field, and the TV
vector.

True $(x,y)$ position is never an input to DQN, RND, or LP-RND. Position and
$(x,y,\mathrm{direction})$ are privileged information used only for logging,
coverage, distractor diagnostics, and heatmaps.

\paragraph{Compared agents.}
\begin{enumerate}
    \item \textbf{DQN}: extrinsic reward only and $\epsilon$-greedy exploration.
    \item \textbf{DQN+RND}: a fixed random target and a trainable predictor,
    with prediction error as intrinsic reward.
    \item \textbf{DQN+LP-RND}: RND gated by progress relative to a lagged
    predictor snapshot.
\end{enumerate}

All methods share the same Double-DQN backbone and hyperparameters. Let $o_t$
denote the complete partial observation, $f$ the frozen RND target,
$g_{\theta}$ the current predictor, and $g_{\bar{\theta}}$ a lagged predictor.
The errors are
\begin{align}
e_{\mathrm{cur}}(o_t) &=
  \lVert g_{\theta}(o_t)-f(o_t)\rVert_2^2,\\
e_{\mathrm{old}}(o_t) &=
  \lVert g_{\bar{\theta}}(o_t)-f(o_t)\rVert_2^2.
\end{align}
Standard RND uses $r_t^{\mathrm{RND}}=e_{\mathrm{cur}}(o_t)$. LP-RND uses
\begin{equation}
r_t^{\mathrm{LP}} =
\max\!\left(0,e_{\mathrm{old}}(o_t)-e_{\mathrm{cur}}(o_t)\right).
\end{equation}
Thus, high but non-reducible prediction error should yield little LP reward.
Training uses $r'_t=r_t^{\mathrm{ext}}+\beta r_t^{\mathrm{int}}$, after online
normalization and clipping of the intrinsic signal. Evaluation uses only
extrinsic reward.

\section{Experiments}
\paragraph{Metrics.}
Task performance comprises greedy evaluation return, success rate, area under
the success-rate curve, and steps to the first and third reward discoveries.
Failure-mode diagnostics comprise position coverage, unique
$(x,y,\mathrm{direction})$ states, hashes of full partial observations,
training fraction in the TV region, visitation heatmaps, and RND/LP signals
inside versus outside the TV region.

\paragraph{Protocol.}
I first run a three-seed pilot with DQN, RND, and LP-RND for 50k steps in noisy
FourRooms, using $\beta=0.01$ for both intrinsic methods. The pilot checks
whether reward is discovered and whether ordinary RND exhibits the expected
noisy-TV failure. I scale up only if this gate is passed.

The main noisy-TV experiment uses five seeds. DQN is run once per seed; RND and
LP-RND use the small pre-specified sensitivity set
$\beta\in\{0.01,0.05,0.1\}$, giving 35 runs in total. The beta sweep is an
ablation, not a broad hyperparameter search. If time remains after the main
experiment, I run a clean FourRooms sanity check using the selected beta.

\section{Expected Outcome and Scope}
I expect standard RND to retain high prediction error at the stochastic
distractor and sometimes sacrifice extrinsic performance. LP-RND should reduce
the inside-TV intrinsic signal and ideally reduce distractor visits while
preserving useful exploration. Count-based exploration, additional
environments, and more complex ensemble or recurrent methods are outside the
main scope.

\paragraph{AI assistance disclosure}
I used ChatGPT as a peer for discussing the idea, challenging design choices,
wording, implementation support, and debugging. I will inspect the code,
execute the experiments, and verify all analysis and final claims myself.

\bibliography{references}
\bibliographystyle{plainnat}

\end{document}
